\documentclass[a4paper,12pt]{article}
\usepackage[utf8]{inputenc}
\usepackage[T2A]{fontenc}
\usepackage[russian]{babel}
\usepackage{amsmath, amssymb, amsthm}
\usepackage{hyperref}
\usepackage{geometry}
\geometry{top=2cm, bottom=2cm, left=2cm, right=2cm}
\newtheorem{definition}{Определение}
\newtheorem{theorem}{Теорема}
\newtheorem{example}{Пример}
\renewcommand{\thesection}{\arabic{section}}


\title{Конспект по теории алгоритмов\\ (по материалам учебников)}
\author{}
\date{}

\begin{document}

\maketitle

\section*{Введение}
В конспекте использованы следующие источники:
\begin{itemize}
    \item \textit{Кормен Т., Лейзерсон Ч., Ривест Р., Штайн К.} Алгоритмы: построение и анализ, 3-е изд. --- М.: И.Д. Вильямс, 2013. (далее --- Кормен)
    \item \textit{Хопкрофт Дж., Мотвани Р., Ульман Дж.} Введение в теорию автоматов, языков и вычислений, 2-е изд. --- М.: И.Д. Вильямс, 2002. (далее --- Хопкрофт)
    \item \textit{Кузнецов О.П., Адельсон-Вельский Г.М.} Дискретная математика для инженера. --- М.: Энергоатомиздат, 1988. (далее --- Кузнецов)
    \item \textit{Охотин А.} Математические основы алгоритмов, осень 2025 г. Лекция 1. (далее --- Охотин)
\end{itemize}

\section{Машины Тьюринга, частично рекурсивные функции, РАМ-машины. Тезис Черча. Неразрешимость проблемы останова}

\subsection{Машина Тьюринга (Хопкрофт, гл. 8; Кузнецов, п. 5.2)}
Машина Тьюринга (МТ) --- математическая модель алгоритма, состоящая из:
\begin{itemize}
    \item \textbf{Ленты}, бесконечной в обе стороны, разбитой на ячейки. В каждой ячейке записан один символ из конечного алфавита \( \Sigma \). Специальный символ \( \lambda \) (или \( B \)) обозначает пустую ячейку.
    \item \textbf{Управляющего устройства} с конечным множеством состояний \( Q \). Выделено начальное состояние \( q_0 \) и заключительное (останавливающее) состояние \( q_f \).
    \item \textbf{Головки}, которая в каждый момент обозревает одну ячейку, может читать символ, записывать новый и сдвигаться влево (\( L \)) или вправо (\( R \)).
\end{itemize}
Работа МТ задаётся \textbf{программой} --- конечным набором команд вида \( q_i a_j \to q_k a_l D \), где \( D \in \{L,R\} \). Если в состоянии \( q_i \) головка читает символ \( a_j \), то она записывает \( a_l \), переходит в состояние \( q_k \) и сдвигается в направлении \( D \).

\textbf{Конфигурация} --- запись вида \( \alpha q_i \beta \), где \( \alpha \) и \( \beta \) --- слова на ленте, \( q_i \) --- текущее состояние, головка обозревает первый символ \( \beta \). Начальная конфигурация: \( q_0 w \), где \( w \) --- входное слово. Если МТ достигает конфигурации с заключительным состоянием, говорят, что она \textbf{останавливается}. Функция \( f \) называется \textbf{вычислимой по Тьюрингу}, если существует МТ, которая по входу \( w \) останавливается с \( f(w) \) на ленте (и не останавливается, если \( f(w) \) не определено).

\subsection{Частично рекурсивные функции (Кузнецов, п. 5.3)}
Класс \textbf{частично рекурсивных функций} строится из простейших функций (константа 0, функция следования \( x' = x+1 \), функции тождества \( I_m^n(x_1,\dots,x_n)=x_m \)) с помощью трёх операторов:
\begin{enumerate}
    \item \textbf{Суперпозиция} \( S_m^n \): \( f(x_1,\dots,x_n)=h(g_1(\dots),\dots,g_m(\dots)) \).
    \item \textbf{Примитивная рекурсия} \( R_n \):
    \[
    \begin{cases}
        f(x_1,\dots,x_n,0)=g(x_1,\dots,x_n),\\
        f(x_1,\dots,x_n,y+1)=h(x_1,\dots,x_n,y,f(x_1,\dots,x_n,y)).
    \end{cases}
    \]
    Функции, получаемые без оператора минимизации, называются \textbf{примитивно-рекурсивными}.
    \item \textbf{\(\mu\)-оператор (оператор минимизации)}: 
    \[
    f(x_1,\dots,x_n)=\mu y\,[P(x_1,\dots,x_n,y)],
    \]
    где \( P \) --- предикат; \( f \) возвращает наименьшее \( y \), для которого \( P \) истинно, и не определено, если такого \( y \) нет. Применение \( \mu \)-оператора к примитивно-рекурсивным предикатам даёт \textbf{частично рекурсивные} функции.
\end{enumerate}
Теорема: \textbf{класс частично рекурсивных функций совпадает с классом функций, вычислимых по Тьюрингу} (Кузнецов, п. 5.4; Хопкрофт, гл. 8).

\subsection{РАМ-машины (Охотин, п. 2.2)}
Модель \textbf{RAM} (Random Access Machine) --- идеализированный компьютер с произвольным доступом к памяти.
\begin{itemize}
    \item Память --- счётное множество ячеек \( x_i \), \( i\in\mathbb{Z} \), каждая хранит целое число.
    \item Программа --- последовательность команд (копирование, арифметика, условный и безусловный переходы, остановка). Адресация: прямая (константа \( n \)), прямая (ячейка \( x_n \)), косвенная (\( x_{x_n} \)).
    \item Входные данные: \( x_0 = n \), \( x_1,\dots,x_n \) --- значения. Выходные данные в том же формате.
\end{itemize}
RAM является более реалистичной моделью для анализа времени работы алгоритмов. Обычно предполагается, что каждая команда выполняется за единицу времени, а числа не превосходят по величине объёма памяти (модель равной стоимости). Существует также модель логарифмической стоимости (log-cost), где операции над числом \( i \) требуют \( \log i \) времени.

\subsection{Тезис Черча (Тьюринга) (Хопкрофт, гл. 8.1.4; Кузнецов, п. 5.4)}
\textbf{Тезис Черча--Тьюринга}: всякая функция, которая может быть эффективно вычислена (алгоритмически), вычислима на машине Тьюринга (и, следовательно, частично рекурсивна). Это не доказуемое утверждение, а соглашение о том, что считать алгоритмом. Оно подтверждается тем, что все известные формализации алгоритмов (МТ, частично рекурсивные функции, RAM, нормальные алгоритмы Маркова и др.) эквивалентны.

\subsection{Неразрешимость проблемы останова (Хопкрофт, гл. 9; Кузнецов, п. 5.4)}
\textbf{Проблема останова (Halting problem)}: по описанию машины Тьюринга \( M \) и входному слову \( w \) определить, остановится ли \( M \) на входе \( w \).

\textbf{Теорема}: проблема останова алгоритмически неразрешима, т.е. не существует машины Тьюринга, которая решала бы эту проблему для всех пар \( (M,w) \).

\textit{Идея доказательства (от противного)}. Предположим, существует машина \( H \), которая по коду \( \langle M,w \rangle \) выдаёт «да», если \( M(w) \) останавливается, и «нет» в противном случае. Построим новую машину \( D \), которая на входе \( \langle M \rangle \) работает так: запускает \( H(\langle M, \langle M \rangle \rangle) \); если \( H \) говорит «да», то \( D \) зацикливается, если «нет» --- останавливается. Тогда, подав на вход \( D \) её собственный код, получим парадокс: \( D(\langle D \rangle) \) останавливается тогда и только тогда, когда не останавливается. Противоречие. Следовательно, такой \( H \) не существует.

\section{Анализ сложности алгоритмов. Классы P и NP. Сводимость. NP-полнота. Теорема Кука–Левина}

\subsection{Асимптотические обозначения (Кормен, гл. 3)}
\begin{itemize}
    \item \( O(g(n)) \) --- верхняя граница: \( 0\le f(n)\le c\,g(n) \) для всех \( n\ge n_0 \).
    \item \( \Omega(g(n)) \) --- нижняя граница: \( 0\le c\,g(n)\le f(n) \) для всех \( n\ge n_0 \).
    \item \( \Theta(g(n)) \) --- точная граница: \( f(n)=O(g(n)) \) и \( f(n)=\Omega(g(n)) \).
    \item \( o(g(n)) \) --- \( f(n)/g(n)\to 0 \); \( \omega(g(n)) \) --- \( f(n)/g(n)\to\infty \).
\end{itemize}
Время работы алгоритма обычно оценивается в наихудшем случае как функция размера входа \( n \).

\subsection{Классы P и NP (Кормен, гл. 34; Хопкрофт, гл. 10)}
\begin{itemize}
    \item \textbf{Класс P} --- множество задач разрешения (языков), для которых существует детерминированный алгоритм, решающий задачу за время \( O(n^k) \) для некоторой константы \( k \) (полиномиальное время).
    \item \textbf{Класс NP} --- множество задач, для которых существует недетерминированный алгоритм, работающий полиномиальное время. Эквивалентно: задача принадлежит NP, если для любого положительного ответа существует «свидетельство» (полиномиальной длины), которое можно проверить за полиномиальное время.
\end{itemize}
Примеры задач из NP: задача выполнимости булевых формул (SAT), задача о клике, задача о вершинном покрытии, задача коммивояжёра (решение «да/нет»). Задачи из P: сортировка, поиск кратчайшего пути, проверка принадлежности регулярному языку.

\subsection{Сводимость (Кормен, гл. 34.3; Хопкрофт, гл. 10.1.5)}
\begin{itemize}
    \item \textbf{Сводимость по Карпу (полиномиальная сводимость)}: язык \( A \) сводится к \( B \) (\( A \le_P B \)), если существует функция \( f \), вычислимая за полиномиальное время, такая что \( x\in A \iff f(x)\in B \).
    \item \textbf{Сводимость по Тьюрингу}: задача \( A \) сводится к \( B \), если существует полиномиальный алгоритм, использующий \( B \) как оракул (т.е. решающий подзадачи из \( B \) за единицу времени).
\end{itemize}

\subsection{NP-полнота (Кормен, гл. 34; Хопкрофт, гл. 10)}
Язык \( L \) называется \textbf{NP-полным}, если:
\begin{enumerate}
    \item \( L\in NP \);
    \item Любой язык из NP сводится к \( L \) за полиномиальное время (сводимость по Карпу).
\end{enumerate}
Свойства:
\begin{itemize}
    \item Если для какого-либо NP-полного языка существует полиномиальный алгоритм, то \( P=NP \).
    \item Если \( L \) NP-полна и \( L \le_P M \), причём \( M\in NP \), то \( M \) также NP-полна.
\end{itemize}

\subsection{Теорема Кука–Левина (Кормен, гл. 34.3; Хопкрофт, гл. 10.2)}
\textbf{Теорема Кука–Левина}: задача выполнимости булевых формул (SAT) является NP-полной. Доказательство: SAT \( \in NP \); любая задача из NP сводится к SAT путём кодирования работы недетерминированной машины Тьюринга булевой формулой.

\section{Классы задач по памяти: L, NL, coNL, PSPACE. Соотношения с временными классами (Хопкрофт, гл. 11)}

\subsection{Определения}
\begin{itemize}
    \item \textbf{L (логрифмическая память)} --- задачи, решаемые детерминированной машиной Тьюринга с использованием \( O(\log n) \) ячеек памяти (на отдельной рабочей ленте; входная лента только для чтения).
    \item \textbf{NL} --- то же, но для недетерминированной машины.
    \item \textbf{coNL} --- класс языков, дополнения которых принадлежат NL.
    \item \textbf{PSPACE} --- задачи, решаемые детерминированной машиной Тьюринга с полиномиальной памятью (\( O(n^k) \)).
\end{itemize}

\subsection{Соотношения между классами}
\begin{itemize}
    \item \( L \subseteq NL \subseteq P \subseteq NP \subseteq PSPACE \).
    \item \( NL \subseteq coNL \) (теорема Иммермана–Шелепчени: NL замкнут относительно дополнения, т.е. \( NL = coNL \)).
    \item \( PSPACE = NPSPACE \) (теорема Сэвича: недетерминированная память может быть промоделирована детерминированной с квадратичными затратами).
    \item \( NL \subseteq PSPACE \) (очевидно).
    \item \( P \subseteq PSPACE \), \( NP \subseteq PSPACE \).
\end{itemize}
Точные включения (например, \( P \subsetneq NP \) или \( NL \subsetneq P \)) неизвестны, но считается, что все эти классы различны.

\section{Вероятностные алгоритмы}
Вероятностная машина Тьюринга – это недетерминированная МТ, в которой при каждом недетерминированном выборе выбор делается случайно с равной вероятностью.
\begin{definition}[Класс $\mathcal{BPP}$]
    Язык $L$ принадлежит $\mathcal{BPP}$ (ошибка с двух сторон), если существует вероятностная МТ, работающая за полиномиальное время, такая что:
    \[
    \Pr[\text{машина допускает } x] \ge \frac{2}{3} \quad \text{для } x \in L,\qquad
    \Pr[\text{машина отвергает } x] \ge \frac{2}{3} \quad \text{для } x \notin L.
    \]
\end{definition}
\begin{definition}[Класс $\mathcal{RP}$]
    Язык $L$ принадлежит $\mathcal{RP}$ (односторонняя ошибка), если существует вероятностная МТ за полиномиальное время, такая что:
    \[
    \Pr[\text{допуск}] \ge \frac{1}{2} \quad \text{для } x \in L,\qquad
    \Pr[\text{отвергание}] = 1 \quad \text{для } x \notin L.
    \]
\end{definition}
\begin{definition}[Класс $\mathcal{ZPP}$]
    Язык $L$ принадлежит $\mathcal{ZPP}$ (нулевая ошибка), если существует вероятностная МТ, которая всегда даёт правильный ответ, а ожидаемое время работы полиномиально (или, эквивалентно, $L$ принадлежит $\mathcal{RP} \cap \mathrm{co}\mathcal{RP}$).
\end{definition}
\paragraph{Соотношения:}
\[
\mathcal{P} \subseteq \mathcal{ZPP} \subseteq \mathcal{RP} \subseteq \mathcal{BPP} \subseteq \mathcal{PSPACE}.
\]
(Источник: [Хопкрофт, гл. 11.4])

\section{Подходы к проектированию алгоритмов}

\subsection{Разделяй и властвуй (Кормен, гл. 4; Охотин, лекция 1)}
Метод: разбить задачу на несколько подзадач меньшего размера, решить их рекурсивно, объединить результаты.
\begin{itemize}
    \item \textbf{Сортировка слиянием (Merge sort)} (Кормен, гл. 2.3; Охотин, п. 4): рекурсивно сортируем две половины массива, затем сливаем их за линейное время. Время работы \( O(n\log n) \).
    \item \textbf{Быстрое умножение Карацубы} (Охотин, п. 3): умножение двух \( n \)-разрядных чисел выполняется за \( O(n^{\log_2 3}) \approx O(n^{1.585}) \) вместо \( O(n^2) \). Идея: \( (a_1\cdot2^{n/2}+a_0)(b_1\cdot2^{n/2}+b_0)=a_1b_1\cdot2^n+(a_1b_0+a_0b_1)\cdot2^{n/2}+a_0b_0 \), причём \( a_1b_0+a_0b_1 = (a_1+a_0)(b_1+b_0)-a_1b_1-a_0b_0 \), что требует только трёх умножений половинных чисел.
\end{itemize}

\subsection{Динамическое программирование (Кормен, гл. 15; Охотин, п. 5,7)}
Метод: разбить задачу на перекрывающиеся подзадачи, решить каждую один раз и запомнить результат (табулирование). Применяется для задач с оптимальной подструктурой.
\begin{itemize}
    \item \textbf{Задача о разделении стержня (rod cutting)} (Кормен, гл. 15.1; Охотин, п. 5): по длине стержня \( n \) и ценам \( p_i \) для кусков длины \( i \) найти максимальную выручку. Рекуррентность: \( T_n = \max_{1\le i\le n}(p_i + T_{n-i}) \). Вычисление за \( O(n^2) \) (псевдополиномиальное).
    \item \textbf{Наибольшая общая подпоследовательность (LCS)} (Кормен, гл. 15.4; Охотин, п. 7): для строк \( u \) и \( v \) длины \( m \) и \( n \) длина LCS вычисляется по формуле:
    \[
    T_{i,j}=
    \begin{cases}
        0, & i=0\ \text{или}\ j=0,\\
        T_{i-1,j-1}+1, & u_i=v_j,\\
        \max(T_{i-1,j}, T_{i,j-1}), & u_i\neq v_j.
    \end{cases}
    \]
    Время \( O(mn) \).
    \item \textbf{Расстояние Левенштейна (редакционное расстояние)} (доп. источник). Для двух последовательностей (строк) \( a_1\ldots a_n \) и \( b_1\ldots b_m \) минимальное количество операций вставки, удаления и замены символа, необходимых для превращения одной строки в другую, вычисляется по формуле:
    \[
    d[i,j] = \min \begin{cases}
        d[i-1,j] + 1 & \text{(удаление } a_i\text{)},\\
        d[i,j-1] + 1 & \text{(вставка } b_j\text{)},\\
        d[i-1,j-1] + \chi(a_i,b_j) & \text{(замена или совпадение)},
    \end{cases}
    \]
    где \( \chi(x,y)=0 \) при \( x=y \), иначе \( 1 \). Базовые случаи: \( d[i,0]=i \), \( d[0,j]=j \). Время работы \( O(mn) \), восстановление ответа --- аналогично LCS. Для вычисления только расстояния достаточно \( O(\min(m,n)) \) памяти.
\end{itemize}

\subsection{Жадная стратегия (Кормен, гл. 16)}
На каждом шаге выбирается локально оптимальное решение в надежде, что оно приведёт к глобальному оптимуму. Применима, когда задача обладает свойством жадного выбора и оптимальной подструктурой.
\begin{itemize}
    \item \textbf{Задача о выборе процессов (activity selection)} (Кормен, гл. 16.1): выбор максимального множества непересекающихся по времени процессов. Жадный алгоритм: всегда выбирать процесс с самым ранним временем окончания.
    \item \textbf{Коды Хаффмана} (Кормен, гл. 16.3): построение оптимального префиксного кода путём многократного слияния двух узлов с наименьшими частотами.
\end{itemize}

\subsection{Другие алгоритмы (Кормен)}
\begin{itemize}
    \item \textbf{Двоичный поиск} (Кормен, упр. 2.3.5): в отсортированном массиве поиск элемента выполняется за \( O(\log n) \) сравнений.
    \item \textbf{Быстрое возведение в степень} (Кормен, гл. 31.6): вычисление \( a^b \bmod m \) за \( O(\log b) \) умножений, используя двоичное разложение показателя.
\end{itemize}

\section{Структуры данных}

\subsection{Двоичные деревья поиска (Кормен, гл. 12)}
\begin{itemize}
    \item \textbf{Определение}: бинарное дерево, в котором для каждого узла \( x \) все ключи в левом поддереве меньше \( x.key \), а в правом --- больше.
    \item \textbf{Операции}: поиск (TREE-SEARCH), минимум, максимум, следующий/предыдущий элемент (TREE-SUCCESSOR), вставка (TREE-INSERT), удаление (TREE-DELETE). Все операции в худшем случае выполняются за \( O(h) \), где \( h \) --- высота дерева.
    \item \textbf{Обходы}: инфиксный (inorder), прямой (preorder), обратный (postorder).
\end{itemize}

\subsection{Кучи (Кормен, гл. 6)}
\begin{itemize}
    \item \textbf{Двоичная куча} --- массив, который можно рассматривать как почти полное бинарное дерево. \textbf{Невозрастающая куча} (max-heap): для каждого узла \( i \) выполняется \( A[PARENT(i)] \ge A[i] \).
    \item \textbf{Поддержка свойства кучи}: процедура MAX-HEAPIFY\( (A,i) \) «проталкивает» элемент вниз за \( O(\log n) \).
    \item \textbf{Построение кучи}: BUILD-MAX-HEAP за \( O(n) \).
    \item \textbf{Пирамидальная сортировка (Heapsort)}: строим кучу, затем многократно извлекаем максимум, переставляя его в конец массива. Время \( O(n\log n) \).
    \item \textbf{Приоритетная очередь}: операции MAXIMUM, EXTRACT-MAX, INSERT, INCREASE-KEY реализуются за \( O(\log n) \).
\end{itemize}

\subsection{Хеш-таблицы (Кормен, гл. 11)}
\begin{itemize}
    \item \textbf{Идея}: отображение ключей в ячейки массива (хеш-таблицы) с помощью хеш-функции \( h:U\to\{0,\dots,m-1\} \).
    \item \textbf{Разрешение коллизий}:
        \begin{itemize}
            \item \textit{Метод цепочек}: каждая ячейка содержит список ключей, попавших в неё. При хорошей хеш-функции среднее время операций \( O(1+\alpha) \), где \( \alpha=n/m \) --- коэффициент заполнения.
            \item \textit{Открытая адресация}: коллизии разрешаются последовательным просмотром ячеек (линейное, квадратичное, двойное хеширование). Требует \( \alpha<1 \).
        \end{itemize}
    \item \textbf{Хеш-функции}: деление (\( h(k)=k\bmod m \)), умножение (\( h(k)=\lfloor m(kA\bmod 1)\rfloor \)), универсальное хеширование (случайный выбор из семейства функций).
\end{itemize}

\end{document}
